\section{Hits against waveforms, same events}\label{sec:tiers}
\sectionstate{partial}{\cref{tab:tiers} generated by \texttt{hits\_daq/make\_tiers\_table.py} from pass-1 det3 (ana1106, both runs) and SPS chamber E with the row's own $v$ (DAG 4646000; \texttt{hits\_daq/DAG\_SCORES\_2026-11-13.md})}
\intent{One table, same events, same reference: H0 the production hit chain, H2c the same forward model fitted to the
hit lists (one time and amplitude per strip), W the waveform fit. The model helps the hits but does not replace the
samples. The other tiers (H1, H2a), the estimator variants, the H2c variant and SPS scans, and the readout requirements
are IIb (D02, D03).}

To separate what the forward model brings from what the samples bring, we fit the same model to the hit lists alone:
one time and one amplitude per strip, as a hit-only front end would record them (H2c). It is compared on the same
planes with the production hit chain (H0) and with the waveform fit (W) in \cref{tab:tiers}. On chamber~A the hit
chain is calibrated against the M3 reference itself: the slope of the hit ladder is mapped to the reference angle by a
linear fit trained on the other half of the events. This is more knowledge than any real hit-only measurement has, and
it still leaves a resolution of 2.1\textdegree{} (X) and 2.6--2.7\textdegree{} (Y), with a quarter (X) to a third (Y) of the
planes more than 5\textdegree{} from the reference. Fitting the forward model to the same hits roughly halves the
spread, to 1.2--1.3\textdegree, with no reference involved; the waveform fit reaches 1.0\textdegree, with 96--97\,\% of
the planes within 5\textdegree. The two runs agree to 0.07\textdegree{} in every column. These widths include the error
of the M3 reference; the chamber's own resolution is in \cref{sec:floor}.

H2c gains because it never reads the angle off the slope of the hit ladder. For a trial track it generates the
model waveforms, runs the same hit-finding on them, and compares the predicted hit list with the observed one, so the
compression of the ladder by charge from neighbouring depths~\cite{paperIIb} is part of the prediction rather than
an error in the measurement. What it lacks is the shape of each pulse: one time and one amplitude per strip
constrain the depth profile of the charge only weakly.

At SPS, chamber~E sat on a stage at 25.4\textdegree{} in Ar/CO$_2$, so the inclined view has a known angle and no
reference track is needed. There is no independent drift velocity on these rows (the charge column outlasts the
64-sample window), so $v$ is set per row from the waveform fit at the stage angle: W's mean is therefore 0 by
construction, and only its spread and tails are tested; its accuracy at a known angle with an independent $v$ is
shown in~\cite{paperI}. With that same $v$ the hit line reads 1.1--1.5\textdegree{} too steep, falling with the
drift field, and H2c removes the bias to within 0.5\textdegree{} but not the spread: 1.5--2.0\textdegree, with
30--34\,\% of the planes outside 5\textdegree, against 1.0--1.1\textdegree{} for W. On these rows about a fifth of
the H2c fits run away from a starting line far off the track. The model therefore helps the hits but does not
replace the samples.

\askDylan{W has only 77--79\,\% of the SPS planes within 5\textdegree{} of the stage angle, against 96--97\,\% on
the bench, while its core is as narrow (1.0--1.1\textdegree). The DAG scores do not say what the other fifth is
(several tracks, showers, the window cut). Is it known from the Paper~I SPS selection, or should it be stated as
not understood?}

\begin{table}[t]
  \centering
  \small
  \caption{Track angle from the hits and from the waveforms on the same planes. Hits: on chamber~A the production
  hit line H0, mapped to the reference angle by a linear fit trained on the other half of the events; on chamber~E the
  hit line with the row's drift velocity. H2c: the forward model fitted to the hit lists (one time and amplitude per
  strip). W: the forward model fitted to the sampled waveforms. $d$ is the reconstructed minus the reference angle;
  $s_{68}$ is the 68th percentile of $|d-\tilde d|$ over $|d|\le5$\textdegree{} ($\tilde d$ the median), in5 the
  fraction of planes with $|d|\le5$\textdegree, and steep the median of $d\,\mathrm{sign}(\tan\theta_\mathrm{ref})$
  over $|d|\le5$\textdegree{} ($>0$: reads too steep). Chamber~A: pass 1, Run 1 and Run 2 of~\cite{paperI}, against
  M3 (the widths include the reference's error). Chamber~E: 3000 planes per row; $v$ is set per row from W at the
  stage angle, so W's steep is 0 by construction.}
  \label{tab:tiers}
  % generated by hits_daq/make_tiers_table.py -- do not edit by hand
\begin{tabular}{@{}lr rrr rrr rrr@{}}
\toprule
 & & \multicolumn{3}{c}{hits} & \multicolumn{3}{c}{H2c} & \multicolumn{3}{c}{W} \\
\cmidrule(lr){3-5}\cmidrule(lr){6-8}\cmidrule(l){9-11}
 & planes & $s_{68}$ & in5 & steep & $s_{68}$ & in5 & steep & $s_{68}$ & in5 & steep \\
 & & [\textdegree] & [\%] & [\textdegree] & [\textdegree] & [\%] & [\textdegree] & [\textdegree] & [\%] & [\textdegree] \\
\midrule
\multicolumn{11}{@{}l}{\emph{Chamber A, cosmic bench, against the M3 reference}} \\
Run 1, X & 6191 & 2.06 & 77 & $+0.35$ & 1.20 & 86 & $-0.05$ & 1.04 & 97 & $-0.01$ \\
Run 1, Y & 6175 & 2.68 & 65 & $+0.34$ & 1.28 & 84 & $+0.19$ & 1.01 & 97 & $+0.03$ \\
Run 2, X & 19202 & 2.11 & 76 & $+0.41$ & 1.20 & 86 & $+0.11$ & 1.04 & 96 & $+0.06$ \\
Run 2, Y & 19162 & 2.61 & 65 & $+0.27$ & 1.29 & 84 & $+0.25$ & 1.01 & 96 & $+0.04$ \\
\midrule
\multicolumn{11}{@{}l}{\emph{Chamber E, SPS, CO$_2$, inclined view at the 25.4\textdegree{} stage, against the stage angle}} \\
run~57, 191\,V/cm & 3000 & 1.47 & 53 & $+1.52$ & 1.99 & 66 & $-0.36$ & 1.10 & 79 & $+0.07$ \\
run~57, 208\,V/cm & 3000 & 1.40 & 55 & $+1.27$ & 1.73 & 68 & $-0.47$ & 1.02 & 78 & $+0.05$ \\
run~57, 226\,V/cm & 3000 & 1.44 & 57 & $+1.22$ & 1.60 & 69 & $-0.37$ & 1.03 & 77 & $+0.06$ \\
run~57, 243\,V/cm & 3000 & 1.45 & 57 & $+1.15$ & 1.57 & 67 & $-0.23$ & 1.03 & 77 & $+0.05$ \\
run~59, 243\,V/cm & 3000 & 1.49 & 60 & $+1.07$ & 1.50 & 70 & $-0.30$ & 1.03 & 77 & $+0.08$ \\
\bottomrule
\end{tabular}

\end{table}

A hit-only front end, and what it would have to record to reach the waveform result, are the subject of the
companion paper~\cite{paperIIb}.
